\documentclass[11pt,letterpaper,oneside]{report}
\usepackage[T1]{fontenc}
\usepackage[utf8]{inputenc}
\usepackage{lmodern}
\usepackage[letterpaper,margin=0.92in,headheight=15pt]{geometry}
\usepackage{amsmath,amssymb,amsthm,mathtools}
\usepackage{microtype}
\usepackage[dvipsnames]{xcolor}
\usepackage{graphicx}
\usepackage{booktabs,longtable,array,tabularx}
\usepackage{enumitem}
\usepackage{fancyhdr}
\usepackage{titlesec}
\usepackage[most]{tcolorbox}
\usepackage{listings}
\usepackage{tikz}
\usetikzlibrary{positioning}
\usepackage{hyperref}
\hypersetup{colorlinks=true,linkcolor=MidnightBlue,urlcolor=RoyalBlue,pdfauthor={Independent Mathematics and Engineering Study Suite},pdftitle={Ordinary Differential Equations Technical Reference}}
\definecolor{ink}{HTML}{162639}
\definecolor{navy}{HTML}{173B61}
\definecolor{blue}{HTML}{2874A6}
\definecolor{teal}{HTML}{007C83}
\definecolor{amber}{HTML}{8B5E16}
\definecolor{pale}{HTML}{EDF4F8}
\definecolor{paleteal}{HTML}{ECF8F5}
\definecolor{paleamber}{HTML}{FFF6E8}
\titleformat{\chapter}[display]{\normalfont\huge\bfseries\color{navy}}{\filleft\fontsize{43}{46}\selectfont\thechapter}{0.4ex}{\titlerule\vspace{1ex}}[\vspace{1ex}\titlerule]
\titleformat{\section}{\normalfont\Large\bfseries\color{navy}}{\thesection}{0.8em}{}
\titleformat{\subsection}{\normalfont\large\bfseries\color{teal}}{\thesubsection}{0.75em}{}
\pagestyle{fancy}\fancyhf{}\fancyhead[L]{\small\color{navy}Differential Equations Engineering Suite}\fancyhead[R]{\small\color{navy}Volume 2 of 8}\fancyfoot[C]{\thepage}\renewcommand{\headrulewidth}{0.3pt}
\setlist[itemize]{leftmargin=*,itemsep=3pt,topsep=5pt}\setlist[enumerate]{leftmargin=*,itemsep=4pt,topsep=5pt}
\setlength{\parskip}{5pt}\setlength{\parindent}{0pt}\setlength{\emergencystretch}{2em}
\newtheorem{theorem}{Theorem}[chapter]
\newtheorem{proposition}[theorem]{Proposition}
\newtheorem{lemma}[theorem]{Lemma}
\theoremstyle{definition}\newtheorem{definition}[theorem]{Definition}
\newtheorem{example}[theorem]{Worked example}
\theoremstyle{remark}\newtheorem{remark}[theorem]{Remark}
\newtcolorbox{bookground}[1][]{enhanced,breakable,colback=pale,colframe=blue,boxrule=0.6pt,title={TEXTBOOK ALIGNMENT},fonttitle=\sffamily\bfseries\small,#1}
\newtcolorbox{modern}[1][]{enhanced,breakable,colback=paleteal,colframe=teal,boxrule=0.6pt,title={MODERN COMPUTATIONAL EXTENSION},fonttitle=\sffamily\bfseries\small,#1}
\newtcolorbox{caution}[1][]{enhanced,breakable,colback=paleamber,colframe=amber,boxrule=0.6pt,title={ASSUMPTION / VALIDATION},fonttitle=\sffamily\bfseries\small,#1}
\newcommand{\dd}{\,\mathrm{d}}
\newcommand{\R}{\mathbb{R}}
\newcommand{\lap}{\mathcal{L}}
\newcommand{\norm}[1]{\left\lVert#1\right\rVert}
\newcommand{\abs}[1]{\left\lvert#1\right\rvert}
\lstdefinestyle{engpy}{language=Python,basicstyle=\ttfamily\footnotesize,keywordstyle=\color{RoyalBlue}\bfseries,commentstyle=\color{gray},stringstyle=\color{teal},backgroundcolor=\color{pale},frame=single,rulecolor=\color{blue!30},breaklines=true,columns=fullflexible,keepspaces=true,showstringspaces=false,tabsize=4,upquote=true}
\lstset{style=engpy}
\begin{document}
\begin{titlepage}
\pagecolor{navy}\color{white}
\vspace*{0.75in}
{\sffamily\bfseries\large MATHEMATICS \& ENGINEERING DOCUMENTATION SUITE\\[12pt]}
{\color{cyan!40}\rule{\textwidth}{2pt}}\\[1.0in]
{\sffamily\bfseries\fontsize{29}{36}\selectfont Ordinary Differential Equations Technical Reference\par}
\vspace{0.45in}
{\Large Volume 2 of 8}\par\vspace{0.55in}
{\color{cyan!30}\rule{0.50\textwidth}{1pt}}\par\vspace{0.3in}
{\large Independent technical companion aligned with the\par}
{\large 13-chapter organization of Nagle, Saff \& Snider (7th ed.)\par}
\vfill
{\large Original derivations, worked examples, verification and engineering practice\par}\vspace{0.3in}
{\small Prepared October 2026 \quad\textbullet\quad pdf\LaTeX{} compatible\par}
\end{titlepage}
\nopagecolor\color{ink}\pagenumbering{roman}
\chapter*{Reader's guide and source boundaries}\addcontentsline{toc}{chapter}{Reader's guide and source boundaries}
This independent study resource is organized around \emph{Fundamentals of Differential Equations and Boundary Value Problems}, Seventh Edition, by R.~Kent Nagle, Edward~B. Saff, and Arthur~David Snider (Pearson, 2018; ISBN 978-0-321-97710-6). The book's chapter and section numbers are given as \emph{navigation references}, not as claims of reproduced text or proofs.

\begin{bookground}
\textbf{Textbook-grounded scope:} topics, terminology, and chapter/section organization corresponding to the cited textbook. All explanations, derivations, exercise prompts, and diagrams in this volume are newly authored and may emphasize different examples. Check the textbook itself for its exact hypotheses, proofs, exercises, or nomenclature in a particular section.
\end{bookground}
\begin{modern}
\textbf{Distinct extension layer:} Python/NumPy/SciPy workflows, engineering QA and reproducibility practices, automated convergence studies, diagnostic plots, and implementation guidance. These are \emph{not} represented as content supplied by the textbook. API calls are examples and should be verified against the installed package version.
\end{modern}
\begin{caution}
\textbf{Mathematical qualification:} formulas require their stated regularity, domain and boundary/initial conditions. Numerical agreement does not replace a proof; computational outputs are not certified error bounds unless explicitly established.
\end{caution}
\tableofcontents\clearpage\pagenumbering{arabic}
\chapter{First-order differential equations: classification and solution}
\begin{bookground}Correspondence: Chapter 2 (Sections 2.1--2.6), with motion and applications in Chapter 3 and solution concepts from Chapter 1. Terminology includes separable, linear, exact, integrating factor, substitutions and transformations.\end{bookground}
\section{Separable equations (\S2.2)}
A first-order ODE $y'=g(t)h(y)$ permits separation when $h(y)\ne0$:
\[\int\frac{\dd y}{h(y)}=\int g(t)\dd t+C.\]
Find equilibrium solutions $h(y_*)=0$ before division. An implicit relation can define multiple branches; the initial condition selects a branch and interval.
\begin{example}[A bounded population model]
Solve $y'=2y(1-y)$, $y(0)=1/4$. Partial fractions give $\log|y/(1-y)|=2t+C$. With $y(0)/(1-y(0))=1/3$,
\[y(t)=\frac{1}{1+3e^{-2t}},\qquad t\in\R.\]
The result stays in $(0,1)$ and satisfies the initial condition. The parameterized logistic model is a mathematical idealization, not a universal growth law.
\end{example}
\section{Linear equations and integrating factors (\S2.3)}
For $y'+p(t)y=q(t)$, define $\mu(t)=\exp(\int p(t)\dd t)$ on an interval of continuity. Then $(\mu y)'=\mu q$ and
\[y(t)=\frac{1}{\mu(t)}\left[C+\int^t\mu(s)q(s)\dd s\right].\]
\begin{example}[Decaying input response]
For $y'+2y=e^{-t}$, $\mu=e^{2t}$, so $y=Ce^{-2t}+e^{-t}$. With $y(0)=0$, $C=-1$; therefore $y=e^{-t}-e^{-2t}$. Direct substitution returns the right side $e^{-t}$.
\end{example}
\section{Exact equations (\S2.4)}
Write $M(x,y)\dd x+N(x,y)\dd y=0$. If $M_y=N_x$ on a suitable simply connected region and the derivatives are continuous, a potential $F$ exists locally with $F_x=M$, $F_y=N$, and solution $F(x,y)=C$. For $M=2xy+1$ and $N=x^2+2y$, $F=x^2y+x+y^2$.
\section{Special integrating factors and transformations (\S2.5--2.6)}
An integrating factor $\mu(x,y)$ converts a nonexact differential form into an exact one; simple candidates depend only on $x$ or only on $y$. For example, if $(M_y-N_x)/N$ depends only on $x$, a possible integrating factor is $\mu(x)=e^{\int (M_y-N_x)/N\dd x}$ where defined. Transformations such as Bernoulli substitution $v=y^{1-n}$ convert $y'+p y=q y^n$ ($n\ne0,1$) to a linear equation. Treat singular and zero solutions separately.
\chapter{Linear second-order equations}
\begin{bookground}Correspondence: Chapter 4 (Sections 4.1--4.10): characteristic/auxiliary equations, superposition, undetermined coefficients, variation of parameters, variable coefficients and vibration models.\end{bookground}
\section{Homogeneous constant-coefficient theory (\S4.2--4.3)}
The characteristic equation for $a y''+b y'+c y=0$ is $ar^2+br+c=0$, with $a\ne0$. Roots $r_1\ne r_2$ give $c_1e^{r_1t}+c_2e^{r_2t}$; a repeated root $r$ gives $(c_1+c_2t)e^{rt}$; conjugate roots $\alpha\pm i\beta$ give $e^{\alpha t}(c_1\cos\beta t+c_2\sin\beta t)$.
\begin{example}[Underdamped oscillator]
For $y''+2y'+5y=0$, the roots are $-1\pm2i$. If $y(0)=1$ and $y'(0)=0$, then $y=e^{-t}(\cos2t+\tfrac12\sin2t)$. Differentiation verifies both initial conditions.
\end{example}
\section{Superposition and nonhomogeneous forcing (\S4.4--4.5)}
For a linear operator $L$, $L(y_h+y_p)=g$ when $Ly_h=0$ and $Ly_p=g$. The \emph{method of undetermined coefficients} matches forcing families closed under differentiation, such as exponential-polynomial-sinusoidal inputs. Resonance requires multiplying the trial by a sufficient power of $t$ to avoid overlap with the homogeneous basis.
\begin{example}[Resonant forcing]
For $y''+y=\cos t$, a trial $A\cos t+B\sin t$ lies in the homogeneous kernel. Taking $y_p=\tfrac12 t\sin t$ gives $y_p''+y_p=\cos t$. The full solution is $y=C_1\cos t+C_2\sin t+\tfrac12t\sin t$.
\end{example}
\section{Variation of parameters and the Wronskian (\S4.6)}
For $y''+p y'+q y=g$ and an independent homogeneous basis $y_1,y_2$, let $W=y_1y_2'-y_1'y_2$. A particular solution is
\[y_p=-y_1\int\frac{y_2 g}{W}\dd t+y_2\int\frac{y_1 g}{W}\dd t.\]
The formula assumes continuity and $W\ne0$ on the chosen interval. If the equation is $a(t)y''+b(t)y'+c(t)y=f(t)$, first divide by $a(t)$ where $a(t)\ne0$ so $g=f/a$.
\section{Variable coefficients and qualitative dynamics (\S4.7--4.8)}
Constant-coefficient characteristic roots do not solve a general variable-coefficient ODE. A Cauchy--Euler equation $t^2y''+a t y'+b y=0$ admits $y=t^m$ for $t>0$; substituting yields $m(m-1)+a m+b=0$. Near nonlinear equilibria, a locally valid linearization may predict nearby behavior, except where zero-real-part eigenvalues require higher-order analysis.
\chapter{Higher-order linear equations}
\begin{bookground}Correspondence: Chapter 6 (Sections 6.1--6.4). This generalizes the second-order constructions from Chapter 4.\end{bookground}
\section{Fundamental sets and the Wronskian}
For a regular $n$th-order linear ODE, an $n$-dimensional homogeneous solution space admits a fundamental set of $n$ independent solutions. The general nonhomogeneous solution is $y_h+y_p$. The $n\times n$ Wronskian of candidate solutions tests linear independence at a regular point under the linear ODE hypotheses; avoid treating a zero Wronskian for arbitrary functions as a universal proof of dependence.
\section{Characteristic polynomials and multiplicity}
For $a_n y^{(n)}+\cdots+a_0y=0$ with constants, roots $r$ of $a_nr^n+\cdots+a_0=0$ generate modes. A root of multiplicity $m$ contributes $e^{rt},te^{rt},\ldots,t^{m-1}e^{rt}$. Conjugate pairs combine into real sine/cosine modes.
\begin{example}[Repeated root in order three]
$(D+1)^3y=0$ has the family $y=e^{-t}(C_0+C_1 t+C_2t^2)$; the three functions are independent and the characteristic root $-1$ has multiplicity three.
\end{example}
\section{Variation of parameters in state form}
An $n$th-order equation can be rewritten $\mathbf x'=A(t)\mathbf x+\mathbf f(t)$. With a fundamental matrix $\Phi(t)$,
\[\mathbf x(t)=\Phi(t)\left(\mathbf c+\int_{t_0}^{t}\Phi(s)^{-1}\mathbf f(s)\dd s\right).\]
This connects higher-order variation of parameters with the matrix viewpoint in Chapter 9.
\chapter{Series solutions and singular points}
\begin{bookground}Correspondence: Chapter 8 (Sections 8.1--8.8): Taylor polynomials, analytic coefficients, Cauchy--Euler equations, Frobenius method and special functions.\end{bookground}
\section{Ordinary-point power series}
At an ordinary point where coefficients are analytic (after normalization) one can seek $y=\sum_{n=0}^\infty a_n(x-x_0)^n$. Substitute into the equation, align powers, and derive a recurrence. A formal recurrence must be accompanied by a convergence discussion when a rigorous solution is required.
\begin{example}[Series from the equation $y'=y$]
Write $y=\sum_{n\ge0} a_n t^n$, substitute, and compare $t^n$: $(n+1)a_{n+1}=a_n$. With $a_0=1$, $a_n=1/n!$, recovering $e^t$.
\end{example}
\section{Regular singular points and Frobenius (\S8.6)}
For $y''+P(x)y'+Q(x)y=0$, $x_0$ is a regular singular point when $(x-x_0)P(x)$ and $(x-x_0)^2Q(x)$ are analytic there. Try $y=(x-x_0)^r\sum_{n=0}^\infty a_n(x-x_0)^n$ with $a_0\ne0$. The lowest-power coefficient gives the \emph{indicial equation}. Repeated roots or integer-separated roots may require logarithmic second solutions.
\section{Special functions and engineering context (\S8.8)}
Bessel, Legendre and other special-function equations occur after separation of PDEs in non-Cartesian geometry. The method is not just a mnemonic: coordinate symmetry determines differential operators and admissible boundary behavior, which in turn determine allowable eigenvalues.
\chapter{Selecting a method: decision table}
\begin{longtable}{p{0.29\textwidth}p{0.28\textwidth}p{0.28\textwidth}}
\toprule Structure recognized & Candidate technique & Critical diagnostic \\ \midrule\endhead
$y'=g(x)h(y)$ & Separate; include equilibrium branches & Zeros of $h$; interval of validity \\
$y'+p(x)y=q(x)$ & Integrating factor & Coefficient continuity \\
$M\dd x+N\dd y=0$ & Exactness / special factor & $M_y=N_x$ and region \\
Constant-coefficient linear ODE & Auxiliary equation, forcing trials & Multiplicities and resonance \\
General forced linear second order & Variation of parameters & Nonzero Wronskian \\
Analytic local coefficients & Power series & Convergence domain \\
Regular singular point & Frobenius method & Indicial roots and log modes \\
Matrix system & Eigenvectors / matrix exponential & Defective eigenvalues \\
Discontinuous or impulsive forcing & Laplace transforms (Vol. 3) & Transform existence and causality \\
\bottomrule
\end{longtable}
\chapter{Modern computational extension: symbolic cross-checks}
\begin{modern}The following computer algebra calls, API usage and residual-based QA are original extensions. Symbolic simplification is a useful check but does not substitute for the hypotheses in Chapters 2, 4 and 6.\end{modern}
\begin{lstlisting}
import sympy as sp
x = sp.symbols('x', real=True)
y = sp.Function('y')
expr = sp.Eq(sp.diff(y(x), x) + 2*y(x), sp.exp(-x))
solution = sp.dsolve(expr, ics={y(0): 0})
residual = sp.simplify(
    sp.diff(solution.rhs, x) + 2*solution.rhs - sp.exp(-x)
)
assert residual == 0
assert sp.simplify(solution.rhs.subs(x, 0)) == 0
print(solution)
\end{lstlisting}
For each computer-generated candidate, evaluate the residual, initial/boundary conditions, singularities, branch constraints and domain. Report methods that \emph{fail} as well as those that succeed; symbolic systems sometimes return implicit or special-function solutions requiring interpretation.
\chapter{Exercises, answer sketches and engineering checks}
\section{Problems}
\begin{enumerate}
\item Solve $y'+3y=6$, $y(0)=0$.
\item Solve $y''-3y'+2y=0$, $y(0)=1$, $y'(0)=0$.
\item Check whether $(2xy+2)\dd x+(x^2+4y)\dd y=0$ is exact; find its potential.
\item Find one particular solution to $y''+4y=\sin2t$.
\item Show the characteristic polynomial for $y'''-y''-y'+y=0$ and identify repeated roots.
\item Explain why multiplying by $t$ is needed for $y''-2y'+y=e^t$.
\item Identify whether $x=0$ is a regular singular point of $x^2y''+xy'+y=0$.
\end{enumerate}
\section{Answers and solution checks}
\begin{enumerate}
\item $y=2(1-e^{-3t})$.
\item $y=2e^t-e^{2t}$.
\item $M_y=2x=N_x$; $F=x^2y+2x+2y^2$.
\item $y_p=-\tfrac14 t\cos 2t$; substitute to verify $y_p''+4y_p=\sin2t$.
\item $r^3-r^2-r+1=(r-1)^2(r+1)$; a root $1$ of multiplicity two.
\item $e^t$ is already in the homogeneous kernel (a double root); in fact take $y_p=\tfrac12 t^2e^t$.
\item Yes; after dividing by $x^2$, $xP(x)=1$ and $x^2Q(x)=1$ are analytic.
\end{enumerate}

\appendix
\chapter{References and reproducibility}
\begin{enumerate}
\item R.~K. Nagle, E.~B. Saff, and A.~D. Snider, \emph{Fundamentals of Differential Equations and Boundary Value Problems}, 7th ed., Pearson, 2018, ISBN 978-0-321-97710-6. Primary source for the \emph{topic map and terminology}; cited locally by chapter and section.
\item SciPy project, \emph{solve\_ivp} reference, \url{https://docs.scipy.org/doc/scipy/reference/generated/scipy.integrate.solve_ivp.html} (modern extension).
\item SciPy project, \emph{solve\_bvp} reference, \url{https://docs.scipy.org/doc/scipy/reference/generated/scipy.integrate.solve_bvp.html} (modern extension).
\item SciPy project, \emph{eigsh} reference, \url{https://docs.scipy.org/doc/scipy/reference/generated/scipy.sparse.linalg.eigsh.html} (modern extension).
\item Python, NumPy, SciPy, and SymPy project documentation; computational checks require recording package versions and tolerances.
\end{enumerate}
\section*{Verification record template}
\begin{tabularx}{\textwidth}{>{\bfseries}p{0.30\textwidth}X}\toprule
Problem identifier & \rule{0pt}{15pt}\ \\ \midrule
ODE/PDE and domain & \\[8pt]
Initial/boundary conditions & \\[8pt]
Units and parameter values & \\[8pt]
Analytic residual & \\ [8pt]
Grid/step refinement & \\ [8pt]
Independent solver comparison & \\ [8pt]
Version and reproducibility notes & \\ [8pt]\bottomrule
\end{tabularx}
\end{document}
